\section*{Capítulo \ref{ch:b1:active-margins} --- Margens ativas}

\begin{solution}{exo:b1:active-margins:1}
(a) $2\pi G\,\Delta\rho\,h = 2\pi\times\num{6.674e-11}\times1700\times3000 =
\qty{2.14e-3}{m.s^{-2}}$: a gravidade seria menor em cerca de
\qty{214}{\milli\gal}. (b) Para \qty{3.5}{km},
$2\pi\times\num{6.674e-11}\times1700\times3500 = \qty{2.50e-3}{m.s^{-2}}$,
\qty{250}{\milli\gal}: o tamanho da anomalia negativa observada, logo a fossa está longe do equilíbrio isostático. (c) A fórmula da laje supõe que a rocha ausente se estende sem limite em todos os lados; uma fossa com apenas dezenas de quilômetros de largura retira muito menos massa ao redor do ponto de medida, de modo que sua água, sozinha, dá menos que \qty{250}{\milli\gal}. O restante do déficit está mais fundo: a placa inteira é dobrada para baixo sob a fossa, e sua \omterm{def:b1:layered-earth:layers}{crosta}, mais leve que o \omterm{def:b1:layered-earth:layers}{manto}, ocupa o lugar do \omterm{def:b1:layered-earth:layers}{manto} em sua base.
\end{solution}

\begin{solution}{exo:b1:active-margins:2}
(a) $\qty{82}{mm.yr^{-1}} = \qty{82}{km.Myr^{-1}}$, logo \qty{820}{km} em
\qty{10}{Myr}. (b) $820\sin\ang{45} = \qty{580}{km}$. (c) O caminho até
\qty{600}{km} é $600/\sin\ang{45} = \qty{849}{km}$, percorrido em
$849/82 = \qty{10.3}{Myr}$.
\end{solution}

\begin{solution}{exo:b1:active-margins:3}
$g = 110/\tan\delta$: \qty{191}{km} para \ang{30}, \qty{110}{km} para \ang{45},
\qty{64}{km} para \ang{60}. O arco acima da laje de \omterm{def:b1:geological-maps:attitude}{mergulho} mais suave fica mais
distante de sua fossa.
\end{solution}

\begin{solution}{exo:b1:active-margins:4}
(a) Convergência: $\qty{2.7e6}{m^2.yr^{-1}}/\qty{5.1e7}{m} =
\qty{0.053}{m.yr^{-1}}$, cerca de \qty{53}{mm.yr^{-1}}. Espalhamento:
$\qty{3.0e6}{m^2.yr^{-1}}/\qty{6.7e7}{m} = \qty{0.045}{m.yr^{-1}}$, cerca de
\qty{45}{mm.yr^{-1}}. Ao longo de um comprimento total menor, as fronteiras de subducção consomem fundo oceânico mais depressa do que as dorsais o produzem, e ainda assim consomem menos no total.
(b) O outro \qty{0.3}{km^2} ou algo assim é absorvido ao longo das outras fronteiras convergentes do modelo, como colisões continentais e zonas convergentes dentro dos oceanos, onde as placas são encurtadas em vez de subduzidas; e as dorsais não são as únicas fronteiras divergentes: riftes continentais e algumas \omterm{def:b1:rifts-ridges-margins:transform}{falhas transformantes} também se abrem um pouco.
\end{solution}

\begin{solution}{exo:b1:active-margins:5}
(a) $\bar x = 350$, $\bar z = 290$; os desvios de $x$ são $-150$, $-50$,
$50$, $150$, e os de $z$ são $-140$, $-40$, $40$, $140$, logo
$S_{xx} = 2(150^2 + 50^2) = \num{50000}$ e
$S_{xz} = 2(150\times140 + 50\times40) = \num{46000}$. Então $s = 0.92$,
$a = 290 - 0.92\times350 = \qty{-32}{km}$, e o \omterm{def:b1:geological-maps:attitude}{mergulho} é
$\arctan 0.92 = \ang{42.6}$. (b) $z = 0$ em $x = 32/0.92 = \qty{34.8}{km}$
a partir da fossa. (c) Pelos pontos extremos a inclinação é $280/300 = 0.933$, um
\omterm{def:b1:geological-maps:attitude}{mergulho} de \ang{43.0}: próximo, mas usa apenas dois dos quatro terremotos.
\end{solution}

\begin{solution}{exo:b1:active-margins:6}
(a) $50/\tan\ang{25} + 60/\tan\ang{45} = 107.2 + 60.0 = \qty{167.2}{km}$.
(b) Uma laje de \ang{45} o tempo todo dá \qty{110}{km}: o segmento raso e suave alarga a lacuna em \qty{57}{km}. (c) O segmento profundo é
$z = 50 + (x - 107.2)$; ele atinge $z = 0$ em $x = \qty{57.2}{km}$, para o lado do continente em relação à fossa, como a linha das Curilas faz a \qty{53}{km}: um segmento raso e suave explica o deslocamento das Curilas. (d) $107.2 + 60/\tan\ang{60} = 107.2 + 34.6 =
\qty{141.9}{km}$.
\end{solution}

\begin{solution}{exo:b1:active-margins:7}
(a) O raio é metade do \omterm{def:b1:geological-maps:attitude}{mergulho}, e $r = R\,\theta$ com $\theta$ em radianos:
\ang{22.5} e $6371\times0.3927 = \qty{2502}{km}$ para \ang{45}; \ang{15} e
\qty{1668}{km} para \ang{30}; \ang{35} e \qty{3892}{km} para \ang{70}.
(b) Os arcos acima de lajes com \omterm{def:b1:geological-maps:attitude}{mergulho} mais suave. (c) A previsão é
$44.9/2 = \ang{22.45}$; o ajuste, \ang{21.9}, é
$0.55/22.45 = \qty{2.4}{\percent}$ menor.
\end{solution}

\begin{solution}{exo:b1:active-margins:8}
(a) $q/k$: $0.035/2.5 = \qty{0.014}{K.m^{-1}} = \qty{14}{K.km^{-1}}$ abaixo
do \omterm{def:b1:active-margins:active-margin}{antearco}, $0.080/2.5 = \qty{32}{K.km^{-1}}$ abaixo do retroarco.
(b) $10 + 14\times30 = \qty{430}{\celsius}$ e
$10 + 32\times30 = \qty{970}{\celsius}$. (c) Parte do \omterm{def:b1:heat-budget:heat-flow}{fluxo de calor} na superfície é produzida pela radioatividade da própria \omterm{def:b1:layered-earth:layers}{crosta}, de modo que menos calor atravessa os níveis mais profundos e o gradiente diminui com a profundidade: uma extrapolação linear do gradiente da superfície superestima a temperatura em profundidade.
\end{solution}

\begin{solution}{exo:b1:active-margins:9}
(a) \qty{32}{mm.yr^{-1}} é \qty{32}{km.Myr^{-1}}; a meia-taxa é
\qty{16}{mm.yr^{-1}}. (b) $100/32 = \qty{3.1}{Myr}$. (c) A taxa não precisa ter sido constante; parte da largura
pode vir do estiramento e da formação de riftes do arco antes que o espalhamento começasse; e o eixo pode ter saltado ou se espalhado de modo assimétrico.
\end{solution}

\begin{solution}{exo:b1:active-margins:10}
(a) O caminho é $600/\sin\ang{45} = \qty{849}{km}$, percorrido a
\qty{82}{km.Myr^{-1}} em \qty{10.3}{Myr}, isto é \qty{3.27e14}{s}.
(b) $\sqrt{\num{e-6}\times\num{3.27e14}} = \qty{1.8e4}{m}$, cerca de
\qty{18}{km}. (c) $(\qty{5e4}{m})^2/\qty{e-6}{m^2.s^{-1}} = \qty{2.5e15}{s}$,
cerca de \qty{79}{Myr}. (d) O calor só alcançou os \qty{18}{km} externos de cada lado da laje: seu núcleo, a \qty{600}{km} de profundidade, continua quase tão frio quanto ao sair da fossa, mais frio e mais resistente que o \omterm{def:b1:layered-earth:layers}{manto} ao redor, o que explica por que ainda pode romper em terremotos nessa profundidade.
\end{solution}

\begin{solution}{exo:b1:active-margins:11}
(a) À distância $L$ atrás da ponta, a superfície fica $L\tan\alpha$ acima do nível da ponta, e a base fica $L\tan\beta$ abaixo dela: o dorso do prisma tem altura $L(\tan\alpha + \tan\beta)$, e o triângulo tem a área
$\frac12 L\times L(\tan\alpha + \tan\beta)$. (b) $\tan\ang{4} + \tan\ang{6} =
0.175$. O prisma recebe $\qty{1}{km}\times\qty{40}{km.Myr^{-1}} =
\qty{40}{km^2.Myr^{-1}}$ de seção, logo $A = 40t$ e
$L = \sqrt{2A/0.175}$: \qty{21.4}{km} após \qty{1}{Myr}, \qty{42.8}{km}
após \qty{4}{Myr}, \qty{67.6}{km} após \qty{10}{Myr}. $L$ cresce como
$\sqrt t$: quatro vezes mais tempo dá o dobro da largura. (c) $A$ é reduzida à metade, e $L$
dividida por $\sqrt2$: \qty{15.1}{km}, \qty{30.2}{km} e \qty{47.8}{km}.
\end{solution}

\begin{solution}{exo:b1:active-margins:12}
(a) $\partial F/\partial a = -2\sum(z_i - a - s x_i) = 0$ dá
$n\bar z = na + sn\bar x$, isto é $\bar z = a + s\bar x$. (b) Cada $z_i$ e
$\bar z$ crescem de \qty{10}{km}, logo os desvios $z_i - \bar z$ e $S_{xz}$
não mudam: a inclinação é a mesma, e $a = \bar z - s\bar x$ cresce de
\qty{10}{km}. (c) Os desvios $x_i - \bar x$ não mudam, logo a inclinação é a
mesma; $\bar x$ diminui de 20, de modo que $a$ aumenta de $20s$. (d) Nenhum dos dois:
o \omterm{def:b1:geological-maps:attitude}{mergulho} depende apenas da inclinação. Erros comuns a todos os pontos deslocam a reta sem girá-la.
\end{solution}

\begin{solution}{pb:b1:active-margins:1}
\textbf{1.} Os oito pontos ficam perto de uma reta com \omterm{def:b1:geological-maps:attitude}{mergulho} de cerca de
\ang{45}, de \qty{96}{km} de profundidade a \qty{155}{km} da fossa até
\qty{573}{km} de profundidade a \qty{609}{km}.

\textbf{2.} $\sum x_i = 3152$, $\sum z_i = 2751$, $\sum x_i^2 = \num{1442910}$,
$\sum x_i z_i = \num{1290548}$; $\bar x = \qty{394.0}{km}$,
$\bar z = \qty{343.875}{km}$.

\textbf{3.} $S_{xx} = 1442910 - 8\times394^2 = \num{201022}$ e
$S_{xz} = 1290548 - 8\times394\times343.875 = \num{206654}$. Então
$s = 206654/201022 = 1.028$ e $a = 343.875 - 1.02802\times394 =
\qty{-61.2}{km}$: $z = -61.2 + 1.028\,x$.

\textbf{4.} $\arctan1.028 = \ang{45.8}$.

\textbf{5.} Em $x = 61.2/1.028 = \qty{59.5}{km}$ para o lado do continente em relação à fossa. Uma laje que deixasse a fossa com esse \omterm{def:b1:geological-maps:attitude}{mergulho} atingiria a superfície ali; ela a atinge na fossa, logo mergulha mais suavemente perto da fossa e se torna mais íngreme abaixo.

\textbf{6.} Os desvios são $-2.1$, $2.0$, $-7.3$, $5.3$, $7.1$, $2.3$,
$-15.0$ e $8.1$ \unit{km}. O evento 7 fica \qty{15}{km} acima da reta, o evento 8
\qty{8}{km} abaixo dela; os outros ficam a cerca de \qty{7}{km} dela, uma fração
da espessura da laje.

\textbf{7.} $\ang{320.9} - \ang{300} = \ang{20.9}$. Perpendicular à
fossa: $87.4\cos\ang{20.9} = \qty{81.6}{mm.yr^{-1}}$; paralela:
$87.4\sin\ang{20.9} = \qty{31.2}{mm.yr^{-1}}$. A convergência é oblíqua.

\textbf{8.} \qty{81.6}{mm.yr^{-1}} é \qty{81.6}{km.Myr^{-1}}: \qty{816}{km} em
\qty{10}{Myr}.

\textbf{9.} O caminho é $632/\sin\ang{45} = \qty{894}{km}$, percorrido em
$894/81.6 = \qty{11.0}{Myr}$.

\textbf{10.} $122 + 11 = \qty{133}{Myr}$.

\textbf{11.} $t = \qty{10.95}{Myr} = \qty{3.46e14}{s}$ e
$\sqrt{\num{e-6}\times\num{3.46e14}} = \qty{1.86e4}{m}$, cerca de \qty{19}{km}
de cada face: muito menos que os \qty{50}{km} até o meio de uma laje de
\qty{100}{km} de espessura, cujo núcleo ainda está frio.

\textbf{12.} $q/k$: \qty{17.2}{K.km^{-1}} abaixo do fundo oceânico antigo,
\qty{14.0}{K.km^{-1}} abaixo do \omterm{def:b1:active-margins:active-margin}{antearco} e \qty{32.0}{K.km^{-1}} abaixo do
retroarco.

\textbf{13.} Sob o \omterm{def:b1:active-margins:active-margin}{antearco} a laje fria fica apenas algumas dezenas de quilômetros abaixo; ela absorve calor da rocha acima dela, de modo que menos calor chega à superfície.

\textbf{14.} $h = \qty{2.60e-3}{m.s^{-2}}/(2\pi\times\num{6.674e-11}\times1700)
= \qty{3647}{m}$, cerca de \qty{3.6}{km}: tanto quanto, ou um pouco mais que, a
profundidade total da fossa abaixo do fundo ao redor, cerca de \qty{3.5}{km}. A
fossa tem massa faltando que nada compensa.

\textbf{15.} A diferença de densidade é $3300 - 1000 = \qty{2300}{kg.m^{-3}}$,
e $2\pi\times\num{6.674e-11}\times2300\times3500 = \qty{3.38e-3}{m.s^{-2}}$,
\qty{338}{\milli\gal}. A fórmula da laje conta a massa ausente até qualquer
distância; a fossa real é dobrada para baixo apenas em uma faixa de dezenas de quilômetros de largura, então seu déficit é menor, cerca de \qty{260}{\milli\gal}, entre as duas estimativas da laje.

\textbf{16.} A resistência da placa dobrada, que mantém a fossa para baixo e a elevação externa para cima, como faria uma placa elástica, e a tração da laje densa que afunda no \omterm{def:b1:layered-earth:layers}{manto}.

\textbf{17.} $-61.2 + 1.028\times164 = \qty{107.4}{km}$.

\textbf{18.} Sim: \qty{107}{km} fica entre 72 e \qty{173}{km}, e apenas cerca de \qty{2}{km} mais fundo que a média, \qty{105}{km}.

\textbf{19.} $x = (100 + 61.2)/1.028 = \qty{156.8}{km}$.

\textbf{20.} $\ang{45.8}/2 = \ang{22.9}$, e
$6371\times22.9\times\pi/180 = \qty{2546}{km}$.

\textbf{21.} O ajuste, \ang{21.9}, é \ang{1.0} menor, cerca de
\qty{4}{\percent}: a casca dobrada descreve bem a curvatura do arco das Curilas.

\textbf{22.} $x = (105 + 61.2)/1.028 = 166.2/1.028 = \qty{161.7}{km}$.

\textbf{23.} $164 - 161.7 = \qty{2.3}{km}$ mais perto da fossa que a frente observada,
$2.3/164 = \qty{1.4}{\percent}$: os vulcões das Curilas se erguem acima de uma laje
um pouco mais funda que a média, como a questão 18 mostrou.

\textbf{24.} Até o quilômetro mais próximo, a reta da laje das Curilas prevê uma
\omterm{def:b1:active-margins:volcanic-arc}{lacuna arco--fossa} de \textbf{\qty[text-series-to-math=true]{162}{km}} para
vulcões acima da laje a \qty{105}{km} (observada: \qty{164}{km}).
\end{solution}
